\section{Síntesis y ampliación}

Esta clase de Overview avanza desde el mecanismo de selección de información de la attention hasta los límites de seguridad de los agent. En apariencia abarca muchos temas, pero en el fondo puede reducirse a tres expansiones de frontera: primero, el Transformer permite que cada token lea directamente el contexto relevante; segundo, el LLM amplía el conjunto de tareas que puede abordar mediante escala, datos y post-training; tercero, el agent añade fuera del modelo memory, tool, state, communication y environment feedback. Cada expansión aporta nuevas capacidades y, al mismo tiempo, desplaza la verificación desde una sola salida hacia el sistema completo.

\subsection{Diez conclusiones centrales}

\begin{importantbox}{Del mecanismo del modelo a la ingeniería de sistemas}
\begin{enumerate}
  \item La attention selecciona el value mediante query--key matching, y la multi-head proporciona varios subespacios de representación;
  \item El Transformer cambia un emparejamiento global de $O(n^2)$ por rutas de información cortas y paralelismo en el entrenamiento;
  \item El LLM sigue basándose en el next-token objective; la instruction following y la safety provienen de entrenamiento adicional y de restricciones del sistema;
  \item Una emergence curve puede reflejar la scale, pero también puede verse amplificada por un metric threshold;
  \item RLHF, DPO, MoE y data quality son diseños de distintos niveles y no pueden reducirse a «un modelo más grande»;
  \item La memory externa, RAG, continual learning y model editing modifican objetos distintos;
  \item CoT, ToT y Socratic decomposition añaden cómputo intermedio, pero deben acompañarse de un verifier;
  \item Un agent es un observation--action--feedback loop, no una sola llamada al modelo;
  \item La dificultad principal de un sistema multi-agent está en el state, el protocol, la verification y la recovery, no en el arranque en paralelo;
  \item El LLM OS es un mapa de componentes; un despliegue real requiere además permission, isolation, budget, fallback y sandbox.
\end{enumerate}
\end{importantbox}

\subsection{Una tabla que articula toda la clase}

\begin{center}
\begin{longtable}{L{0.17\textwidth} L{0.22\textwidth} L{0.24\textwidth} L{0.25\textwidth}}
\toprule
Nivel & Capacidad nueva & Modo de falla nuevo & Verificación necesaria \\
\midrule
\endhead
Embedding & Representación continua compartida & La similitud no equivale a una relación factual & Pruebas de vecinos cercanos, sesgo y distribución \\
Attention & Lectura del contexto según su contenido & Pesos difíciles de interpretar directamente; costo cuadrático en secuencias largas & ablation, causal test, pruebas de eficiencia \\
LLM scale & Mayor cobertura de tareas y few-shot behavior & Costo, contaminación, capacidades impredecibles & benchmark multidimensional, auditoría de datos \\
Preference tuning & Mayor apego a las preferencias humanas u organizacionales & reward hacking, sesgo de preferencias & held-out preference, evaluación de seguridad \\
External memory & Hechos actualizables y personalización & Recuperación errónea, acceso sin autorización, memory caducada & Verificación de la recuperación, la evidencia y los permisos \\
Reasoning scaffold & Más test-time compute y pasos visibles & Coherente pero erróneo, costo de búsqueda & Herramientas, verificación y pruebas contrafactuales \\
Agent loop & Ejecución de tareas externas de largo plazo & side effect, loop, plan divergence & Bitácora de estado, presupuesto, reversión y recuperación \\
Multi-agent & Paralelismo y specialization & Ambigüedad de protocolo, conflictos, efectos secundarios duplicados & Mensajes estructurados, control de versiones y revisión independiente \\
\bottomrule
\end{longtable}
\end{center}

\subsection{Las preguntas que realmente deja el cierre de la clase}

La closing discussion del ponente concentra el futuro de los agent en error correction, security, permission y sandbox. Quedan además varias preguntas abiertas que vale la pena seguir:

\begin{enumerate}
  \item ¿Cómo medir el long-horizon success, en lugar de medir solo un tool call de un paso?
  \item ¿Cómo lograr personalization sin almacenar indefinidamente datos privados?
  \item ¿Cómo hacer que la memory admita a la vez tiempo, jerarquía, relaciones y olvido?
  \item ¿Cómo distinguir el valor explicativo de una reasoning trace del mecanismo causal real?
  \item ¿Cómo lograr que la multi-agent verification sea independiente del canal que originó el error?
  \item ¿Cómo definir estrategias de ejecución distintas para las action revocables, compensables e irreversibles?
  \item ¿Cómo hacer que los modelos más pequeños, los modelos especializados y la retrieval colaboren en el equilibrio entre capacidad y costo?
\end{enumerate}

\begin{warningbox}{Juicio de ingeniería final}
La action proposal que entrega el modelo debe separarse del side effect real. Todo sistema que involucre dinero, identidad, publicación, eliminación, permisos u operaciones irreversibles debe colocar el modelo detrás de un control plane restringido; una autonomy sin confirmation, audit, recovery ni sandbox no debe presentarse como automatización confiable.
\end{warningbox}

\subsection{Preguntas de autoevaluación}

\begin{enumerate}
  \item ¿Por qué el scaling por $\sqrt{d_k}$ de la attention ayuda a estabilizar el softmax?
  \item ¿En qué se diferencian las fuentes de Q/K/V en la self-attention y en la cross-attention?
  \item ¿Por qué una emergent metric puede presentar un salto aunque la capacidad subyacente varíe de forma suave?
  \item ¿En qué paso utilizan RLHF y DPO, respectivamente, la preference data?
  \item ¿Por qué difieren el número total de parámetros de un MoE y sus active parameters?
  \item ¿Qué modifica cada uno de RAG, fine-tuning, continual learning y model editing?
  \item ¿Por qué los pasos intermedios legibles de CoT no son necesariamente una explicación fiel?
  \item ¿Qué agent component le faltan a una sola LLM call?
  \item ¿Qué riesgos adicionales tiene el UI control frente al API control?
  \item ¿Por qué el manager no puede confiar directamente en el ``done'' del worker?
  \item Cuando una action no puede repetirse, ¿qué mecanismo debe añadirse al redo loop?
  \item En la analogía del LLM OS, ¿qué garantías de un OS tradicional siguen faltando?
\end{enumerate}

\subsection{Lecturas adicionales}

\begin{itemize}
  \item \href{https://arxiv.org/abs/1706.03762}{Vaswani et al., Attention Is All You Need}: el artículo original del Transformer y la multi-head attention.
  \item \href{https://arxiv.org/abs/2206.07682}{Wei et al., Emergent Abilities of Large Language Models} y \href{https://arxiv.org/abs/2304.15004}{Schaeffer et al., Are Emergent Abilities of Large Language Models a Mirage?}: para entender, respectivamente, el fenómeno y la metric caveat.
  \item \href{https://arxiv.org/abs/2305.18290}{Rafailov et al., Direct Preference Optimization}: el DPO objective y su relación con RLHF.
  \item \href{https://arxiv.org/abs/2202.05262}{Meng et al., Locating and Editing Factual Associations in GPT}: ROME y la factual model editing.
  \item \href{https://arxiv.org/abs/2201.11903}{Wei et al., Chain-of-Thought Prompting} y \href{https://arxiv.org/abs/2305.10601}{Yao et al., Tree of Thoughts}: la rationale lineal y el reasoning basado en búsqueda.
  \item \href{https://docs.google.com/presentation/d/1oXPs3LXtIVIsVbwTyGjAWj_aWvak9c1uNC4uhkS6glk/edit}{CS25 V4 official slide deck}: la canonical source de 114 diapositivas de esta clase.
\end{itemize}

\end{document}
